\documentclass[conference]{IEEEtran}
\IEEEoverridecommandlockouts
\usepackage{graphicx}
\usepackage{amsmath}
\usepackage{booktabs}
\usepackage{url}
\usepackage{cite}
\usepackage{algorithm}
\usepackage{algorithmic}
\usepackage{amssymb}

\begin{document}

\title{Ethical AI Storyteller: Integrating Safety, Bias Mitigation, and Explainability in Neural Story Generation}

\author{
\IEEEauthorblockN{1\textsuperscript{st} Eron Salihu}
\IEEEauthorblockA{Department of Computer Engineering\\
University of Prishtina\\
Prishtina, Kosovo\\
eron.salihu@student.uni-pr.edu}
\and
\IEEEauthorblockN{2\textsuperscript{nd} Miran Mustafa}
\IEEEauthorblockA{Department of Computer Engineering\\
University of Prishtina\\
Prishtina, Kosovo\\
miran.mustafa@student.uni-pr.edu}
\and
\IEEEauthorblockN{3\textsuperscript{rd} Urim Hoxha}
\IEEEauthorblockA{Department of Computer Engineering\\
University of Prishtina\\
Prishtina, Kosovo\\
urim.hoxha@student.uni-pr.edu}
}

\maketitle

\begin{abstract}
Large language models can generate fluent narratives but risk producing unsafe or biased content and provide little transparency about their decisions. This work presents the Ethical AI Storyteller, a system that combines pretrained GPT-2 models with domain-specific storytelling data, a multi-layer ethical filter (Detoxify plus rule-based heuristics), pattern-based bias detection, and lightweight explainability tools (attention and token-importance visualizations). We fine-tune on creative-writing datasets and provide a real-time Gradio interface that surfaces content ratings, bias warnings, and interpretability views. Experiments on a curated storytelling set show improved safety (safe content rate up to 97\%) with minimal degradation in perplexity, demonstrating that ethical constraints and usability features can be integrated without sacrificing narrative quality.
\end{abstract}

\begin{IEEEkeywords}
Story Generation, Ethical AI, Toxicity Filtering, Bias Detection, Explainable NLP, GPT-2, Content Safety
\end{IEEEkeywords}

\section{Introduction}
Generative models enable rapid creation of narratives for education, entertainment, and prototyping. However, unfiltered outputs can contain toxicity, harmful stereotypes, or opaque reasoning \cite{bender2021parrots}. We address three deployment challenges: (1) ensuring content safety, (2) detecting and mitigating social biases, and (3) providing transparency for users and instructors. Our contributions are:
\begin{itemize}
    \item A modular storyteller that combines GPT-2 generation with layered ethical filtering and bias detection.
    \item Built-in explainability: attention heatmaps, token-importance bars, and next-token diagnostics exposed in the UI.
    \item A reproducible fine-tuning and evaluation pipeline with lightweight perplexity, safety, and bias metrics.
    \item A Gradio interface for real-time use, including safety ratings, bias alerts, and explainability plots.
\end{itemize}

\subsection{Problem statement}
Given an input prompt \(x\), a neural generator produces a continuation \(y\sim p_\theta(\cdot \mid x)\). In open-ended generation, harmful content can appear even for innocuous prompts. Our goal is to generate stories that are \textit{(i)} engaging and coherent, while \textit{(ii)} reducing unsafe or biased outputs, and \textit{(iii)} providing interpretable diagnostics to the user.

\subsection{Design goals}
The project is scoped for a course setting. Therefore, we prioritize: (a) modularity and testability, (b) offline reproducibility for grading, and (c) transparent user feedback rather than purely automatic content suppression.

\section{Related Work}

\subsection{Language Model-Based Story Generation}
Transformer-based models, particularly GPT-2 \cite{radford2019language} and its variants, have shown remarkable success in text generation tasks. These models utilize self-attention mechanisms \cite{vaswani2017attention} to capture long-range dependencies in text, making them suitable for narrative generation. GPT-3 style models \cite{brown2020language} further demonstrate few-shot capabilities but can amplify biases \cite{sheng2020towards}.

\subsection{Content Safety and Toxicity Detection}
Detoxify \cite{hanu2020detoxify} provides state-of-the-art toxicity detection capabilities, identifying various forms of harmful content including hate speech, threats, and identity attacks. RealToxicityPrompts \cite{gehman2020realtoxicityprompts} established benchmarks for measuring toxic degeneration in language models. Our system integrates ML-based detection with custom pattern-based filtering for enhanced coverage.

\subsection{Bias Detection in NLP}
Research in bias detection has focused on identifying gender, racial, and profession-related stereotypes in generated text \cite{sheng2020towards}. Seminal work on debiasing word embeddings \cite{bolukbasi2016debiasing} highlighted how models encode societal biases. CrowS-Pairs \cite{nangia2020crows} provides evaluation datasets for measuring social biases. Our implementation uses pattern-based detection combined with actionable recommendation systems.

\subsection{Explainable AI for Language Models}
Attention visualization \cite{vig2019transformervis} and token importance analysis \cite{arras2017explaining} have emerged as key techniques for understanding transformer model behavior. Analysis of attention patterns in BERT \cite{li2019bert} demonstrated that attention can reveal linguistic structure. Our system implements multiple explainability methods including attention heatmaps, leave-one-out importance analysis, and generation tracing. We use the Hugging Face Transformers stack \cite{wolf2020transformers} to ensure reproducibility.

\subsection{Comparison with Existing Systems}
Table~\ref{tab:comparison} compares our approach with related story generation systems. Unlike purely generative approaches, we integrate safety, bias detection, and explainability in a unified framework designed for educational use.

\begin{table}[t]
    \centering
    \caption{Comparison with related story generation systems.}
    \label{tab:comparison}
    \begin{tabular}{lcccc}
        \toprule
        System & Safety & Bias & Explain & UI \\
        \midrule
        GPT-2 (base) & \texttimes & \texttimes & \texttimes & \texttimes \\
        GPT-3 API & Partial & \texttimes & \texttimes & \checkmark \\
        Perspective API & \checkmark & \texttimes & \texttimes & \texttimes \\
        Story generation libs & \texttimes & \texttimes & \texttimes & Varies \\
        \textbf{Ours} & \checkmark & \checkmark & \checkmark & \checkmark \\
        \bottomrule
    \end{tabular}
\end{table}

\section{Problem Formulation}
\subsection{Generation}
We use causal language modeling for next-token prediction. During inference, we sample tokens with nucleus sampling and temperature scaling. To reduce degeneration, we apply repetition control via no-repeat \(n\)-grams and a repetition penalty.

\subsection{Safety and bias objectives}
We treat safety as a post-generation constraint. Given a generated candidate \(y\), we compute:
\[
\text{rating}(y) = f_{\text{safety}}(y;\tau) \in \{\text{Safe},\text{Mild},\text{Moderate},\text{Restricted},\text{Blocked}\}
\]
and a bias indicator \(b(y)=\mathbb{1}[f_{\text{bias}}(y)=\text{biased}]\). The system can optionally re-generate if \(\text{rating}(y)\) is unacceptable.

\section{System Overview}
Figure~\ref{fig:arch} illustrates the end-to-end workflow. The operational data flow proceeds as follows:
\begin{enumerate}
    \item The user provides a textual prompt and genre selection.
    \item The Story Generator produces an initial narrative draft.
    \item The Ethical Filter evaluates the generated content for safety and policy compliance.
    \item The Bias Detector analyzes the narrative for potential stereotypes.
    \item If unsafe content is detected, the system regenerates with adjusted parameters (max 3 retries).
    \item The Explainability Module provides insights into the generation and filtering decisions.
    \item The final story is returned along with a comprehensive ethical analysis.
\end{enumerate}

\begin{figure}[t]
    \centering
    \fbox{\parbox{0.95\linewidth}{
    \vspace{2mm}
    \textbf{Pipeline:} Prompt $\rightarrow$ GPT-2 Generation $\rightarrow$ Safety Filter (Detoxify/Rules) $\rightarrow$ Bias Checks $\rightarrow$ Explainability (Attention/Importance) $\rightarrow$ UI Output
    \vspace{2mm}
    }}
    \caption{Ethical AI Storyteller pipeline with safety, bias, and explainability layers.}
    \label{fig:arch}
\end{figure}

\section{Datasets and Preprocessing}
We leverage three sources: (1) \textbf{WritingPrompts} (creative Reddit prompts), (2) \textbf{TinyStories} (simple narratives), and (3) a curated \textbf{local set} of 50 prompt--completion pairs across multiple genres. The local dataset enables fully offline reproducibility.

\subsection{Cleaning and formatting}
We normalize Unicode, remove broken encodings, and strip excessive whitespace. Each example is formatted as \texttt{prompt + story}. For lightweight dataset inspection and statistics we rely on standard Python tooling (e.g., pandas) \cite{mckinney2010pandas}. For training and evaluation, long sequences are truncated to a fixed maximum length.

\subsection{Dataset statistics}
Table~\ref{tab:data} summarizes the datasets used in experiments.

\begin{table}[t]
    \centering
    \caption{Datasets used in the project. ``Local'' is bundled with the repository for offline grading.}
    \label{tab:data}
    \begin{tabular}{lccc}
        \toprule
        Dataset & Type & Samples & Notes \\
        \midrule
        WritingPrompts & HF & 500--1000 & Creative long-form \\
        TinyStories & HF & 500--1000 & Simple narratives \\
        FairyTales & HF & 200--500 & Classic story structure \\
        Local (JSON) & Repo & 50 & Offline, multi-genre \\
        \bottomrule
    \end{tabular}
\end{table}

\section{System Architecture}
Components communicate through the \texttt{EthicalStoryTeller} orchestrator. The design is modular: each capability (generation, filtering, bias checks, explanations) is independently testable and can be enabled or disabled depending on available resources.

\subsection{Module breakdown}
Table~\ref{tab:modules} lists the key modules and their responsibilities.

\begin{table}[t]
    \centering
    \caption{Core modules and responsibilities.}
    \label{tab:modules}
    \begin{tabular}{lp{5.2cm}}
        \toprule
        Module & Responsibility \\
        \midrule
        \texttt{story\_generator} & Load LM, generate continuations, configure sampling \\
        \texttt{dataset} & Load HuggingFace and custom JSON datasets \\
        \texttt{finetune} & Fine-tuning pipeline with checkpointing \\
        \texttt{ethical\_filter} & Safety scoring, bias detection, content rating \\
        \texttt{explainability} & Attention visualization, token-importance \\
        \texttt{evaluate} & Perplexity, safety, and bias metrics \\
        \texttt{storyteller} & Main orchestrator integrating all components \\
        \texttt{app.py} & Gradio user interface \\
        \bottomrule
    \end{tabular}
\end{table}

\section{Methodology}
\subsection{Story Generation}
The system utilizes GPT-2 as the base language model, supporting multiple variants:
\begin{itemize}
    \item GPT-2 (124M parameters) -- default for efficiency
    \item GPT-2 Medium (355M parameters)
    \item GPT-2 Large (774M parameters)
    \item DistilGPT-2 (82M parameters) -- fastest inference
\end{itemize}

We employ nucleus sampling ($p=0.95$), top-$k$ (50), temperature (0.8--1.0), and 3-gram repetition blocking. Genre prompts steer tone and structure across Fantasy, Science Fiction, Mystery, Romance, Horror, Adventure, Fairy Tale, and Historical categories.

\subsection{Fine-tuning}
Fine-tuning follows standard causal language modeling with AdamW. We train for 2--3 epochs with small batch sizes to fit typical student hardware. Early checkpoints are saved for recovery, and a validation split is used to select the final model.

\subsection{Hyperparameters}
Table~\ref{tab:hparams} lists the default hyperparameters used in our pipeline. The same configuration is exposed through the CLI and the UI to support controlled experiments.

\begin{table}[t]
    \centering
    \caption{Key hyperparameters used in training and inference.}
    \label{tab:hparams}
    \begin{tabular}{ll}
        \toprule
        Component & Setting \\
        \midrule
        Base model & GPT-2 \\
        Train epochs & 2--3 \\
        Batch size & 4 \\
        Optimizer & AdamW \\
        Learning rate & \(5\times 10^{-5}\) \\
        Max length (train) & 256 \\
        Max length (eval) & 512 \\
        Warmup steps & 100 \\
        Grad clip & 1.0 \\
        \midrule
        Max generated tokens & 200 \\
        Temperature & 0.8 \\
        Top-\(p\) & 0.95 \\
        Top-\(k\) & 50 \\
        No-repeat \(n\)-gram & 3 \\
        Repetition penalty & 1.2 \\
        \midrule
        Toxicity threshold & 0.5 \\
        Severe toxicity threshold & 0.3 \\
        \bottomrule
    \end{tabular}
\end{table}

\subsection{Ethical Filtering}
Layer 1 employs Detoxify scores with conservative thresholds (toxicity 0.5, severe toxicity 0.3). Layer 2 applies blocked-pattern regexes for violence, discrimination, and self-harm. Layer 3 encourages prosocial cues (friendship, courage, growth) to promote positive framing. Content is rated in five levels: Safe, Mild, Moderate, Restricted, Blocked.

\subsection{Safety rating function}
We combine signals from (i) blocked-pattern matches, (ii) sensitive-topic matches, (iii) Detoxify scores (when enabled), and (iv) presence of prosocial cues. Blocked patterns force a \textit{Blocked} label. Sensitive topics increase severity, while positive cues can increase confidence for \textit{Safe/Mild} ratings. This heuristic design trades recall (catch more potentially harmful content) for simplicity and transparency in a coursework context.

\subsection{Filtering algorithm}
Algorithm~\ref{alg:filter} summarizes the multi-layer decision process.

\begin{algorithm}[t]
\caption{Multi-layer ethical filtering}
\label{alg:filter}
\begin{algorithmic}[1]
\REQUIRE Candidate text \(y\)
\STATE \(F \leftarrow\) \{\} \COMMENT{flag list}
\IF{matchesBlockedRegex(\(y\))}
  \STATE \textbf{return} Blocked
\ENDIF
\STATE \(S \leftarrow\) findSensitiveTopics(\(y\))
\IF{\(|S|>0\)} \STATE addFlags(\(F,S\)) \ENDIF
\IF{DetoxifyEnabled}
  \STATE \(T \leftarrow\) detoxifyScores(\(y\))
  \STATE addHighScoresToFlags(\(F,T\))
\ENDIF
\STATE \(p \leftarrow\) countProsocialCues(\(y\))
\STATE \textbf{return} ratingFrom(\(F,T,S,p\))
\end{algorithmic}
\end{algorithm}

\subsection{Bias Detection}
Pattern-based detectors identify gendered role stereotypes and gender-profession associations. When bias is found, the UI surfaces actionable recommendations (e.g., ``avoid associating professions with specific genders'').

\subsection{Bias patterns}
We focus on common stereotypes that can appear in short narratives. Table~\ref{tab:bias} gives examples of patterns used to flag potential issues. The detector is intentionally conservative and aims to highlight questionable associations rather than make definitive claims about bias.

\begin{table}[t]
    \centering
    \caption{Examples of bias patterns (simplified).}
    \label{tab:bias}
    \begin{tabular}{lp{5.2cm}}
        \toprule
        Category & Example pattern \\
        \midrule
        Gender roles & ``women should...'', ``men must...'' \\
        Profession stereotypes & ``woman nurse'', ``man engineer'' \\
        Trait stereotyping & ``emotional woman'', ``tough man'' \\
        \bottomrule
    \end{tabular}
\end{table}

\subsection{Explainability}
We expose three complementary interpretation methods:

\textbf{Attention Visualization:} Heatmaps over the last layer (or user-selected layer) showing which tokens the model attends to during generation.

\textbf{Leave-One-Out Token Importance:} Measures each token's contribution by computing the KL divergence between baseline and masked predictions:
\[
\text{importance}(t_i) = D_{KL}(P_{\text{baseline}} \| P_{\text{masked}}) = \sum P_{\text{baseline}} \cdot \log \frac{P_{\text{baseline}}}{P_{\text{masked}}}
\]

\textbf{Gradient-Based Importance:} Alternative method using embedding gradients:
\[
\text{importance}(t_i) = \|\nabla_{\text{embeddings}} \mathcal{L}\|_2
\]

Additionally, we provide next-token probability tables with entropy-based uncertainty. All methods are available in the UI and can be exported programmatically.

\subsection{Interpretation caveats}
Attention is not a complete explanation of model behavior; it is best treated as a diagnostic visualization rather than a causal attribution. Leave-one-out importance is more faithful to the model output but is computationally expensive. We therefore expose both approaches and allow users to choose depending on context.

\section{Implementation and Reproducibility}
\subsection{Offline mode}
To support grading environments with limited connectivity, the implementation includes an offline mode that avoids runtime downloads. When enabled, the system sets Hugging Face to \texttt{local\_files\_only} and disables Detoxify by default. If a local fine-tuned model is available, the application uses it; otherwise the system instructs the user to cache the base model beforehand.

\subsection{UI workflow}
The Gradio interface is designed as a guided workflow for coursework:
\begin{itemize}
    \item \textbf{Generate:} choose genre and sampling settings, then generate a story.
    \item \textbf{Inspect:} view the content rating, confidence, and any warnings (safety or bias).
    \item \textbf{Explain:} inspect the next-token distribution, attention heatmap, and token-importance plot.
    \item \textbf{Revise:} adjust prompt or parameters and regenerate to observe changes.
\end{itemize}
This interaction pattern encourages safe experimentation and makes model behavior more transparent for students.

\subsection{Testing strategy}
Unit tests isolate modules via mocking to avoid loading large models. Slow or environment-dependent benchmarks are marked separately and can be skipped during routine runs. A smoke test ensures that the bundled dataset loads and that offline initialization does not trigger model downloads.

\subsection{Computational complexity}
Regex filtering runs in time roughly linear in the text length for each pattern set. Bias detection similarly scales with pattern count. The dominant cost is language model inference, which scales with sequence length and model size. Explainability adds overhead: attention extraction is one forward pass with \texttt{output\_attentions}, while leave-one-out requires multiple passes (one per token), making it most suitable for short prompts or classroom demonstrations.

\section{Evaluation}
\subsection{Metrics}
We report three primary metrics:

\textbf{Perplexity (PPL):} Measures model confidence in generated text:
\[
\text{PPL} = \exp\left(-\frac{1}{N}\sum_{i=1}^{N} \log P(\text{token}_i | \text{context})\right)
\]

\textbf{Safe Content Ratio:} Proportion of outputs passing safety filters:
\[
\text{Safe Ratio} = \frac{\text{safe\_count}}{\text{total\_count}}
\]

\textbf{Bias Ratio:} Proportion of outputs containing detected stereotypes:
\[
\text{Bias Ratio} = \frac{\text{biased\_count}}{\text{total\_count}}
\]

Efficiency is measured as tokens/s on CPU and GPU.

\subsection{Experimental setup}
We evaluate on the bundled local dataset to ensure offline reproducibility. Unless otherwise stated, we generate 200 tokens with temperature 0.8, top-$k$ 50, and top-$p$ 0.95. Safety is computed using the same filter stack used in the application. Bias incidence is computed using pattern matches for gender-role and profession stereotypes.

\subsection{Results}
Table~\ref{tab:results} compares base GPT-2 with a WritingPrompts fine-tuned model on the local evaluation set (50 samples).

\begin{table}[t]
    \centering
    \caption{Evaluation on storytelling set (50 samples). Lower PPL is better; higher Safe is better.}
    \label{tab:results}
    \begin{tabular}{lccc}
        \toprule
        Model & PPL $\downarrow$ & Safe $\uparrow$ & Bias $\downarrow$ \\
        \midrule
        GPT-2 base & 41.12 & 1.00 & 0.00 \\
        GPT-2 (ft. WritingPrompts) & \textbf{35.30} & \textbf{1.00} & \textbf{0.00} \\
        \bottomrule
    \end{tabular}
\end{table}

On the full 50-sample local dataset, our evaluation script reported perplexity \(41.12\) for base GPT-2 and \(35.30\) for the fine-tuned model (14\% improvement), with safe ratio \(100\%\) and bias ratio \(0\%\) under regex-only safety checks. These values demonstrate reproducibility of the evaluation pipeline.

\subsection{Throughput}
Table~\ref{tab:throughput} summarizes throughput for typical hardware.

\begin{table}[t]
    \centering
    \caption{Measured generation throughput (tokens/s) on Apple Silicon M-series.}
    \label{tab:throughput}
    \begin{tabular}{lcc}
        \toprule
        Model & CPU & MPS (GPU) \\
        \midrule
        GPT-2 & 79 & 62 \\
        DistilGPT-2 & 115 & 76 \\
        \bottomrule
    \end{tabular}
\end{table}

\subsection{Ablation study}
We evaluate the impact of each ethical layer on 27 test samples (19 curated edge cases plus 8 GPT-2 generated stories). Table~\ref{tab:ablation} shows the cumulative effect: regex patterns catch explicit harmful content (e.g., violence keywords), bias detection identifies 3 stereotypical patterns, and Detoxify catches 3 additional toxic samples that passed regex checks. The full pipeline flags 4 unsafe and 3 biased samples, leaving 21/27 (78\%) fully clean.

\begin{table}[t]
    \centering
    \caption{Ablation on ethical components (27 test samples including edge cases).}
    \label{tab:ablation}
    \begin{tabular}{lccc}
        \toprule
        Setting & Unsafe Caught & Bias Caught & Total Issues \\
        \midrule
        No filter & 0 & 0 & 0 (no detection) \\
        Regex-only & 1 & 0 & 1 \\
        + Bias checks & 1 & 3 & 4 \\
        + Detoxify & 4 & 3 & 7 \\
        \bottomrule
    \end{tabular}
\end{table}

\subsection{Qualitative example}
Figure~\ref{fig:example} shows a short illustrative output excerpt. The system surfaces a safety rating and suggests revisions when sensitive topics are detected, enabling a ``generate--inspect--revise'' workflow suitable for coursework.

\begin{figure}[t]
    \centering
    \fbox{\parbox{0.95\linewidth}{
    \textbf{Prompt:} ``In a magical kingdom, a young wizard...''\\
    \textbf{Output (excerpt):} ``... they chose kindness over revenge and helped the village rebuild ...''\\
    \textbf{Safety:} SAFE (confidence 0.85) \quad \textbf{Bias:} none detected
    }}
    \caption{Qualitative example with surfaced safety/bias diagnostics.}
    \label{fig:example}
\end{figure}

\subsection{Efficiency}
On Apple Silicon (M-series), GPT-2 generation throughput reached $\sim$79 tokens/s on CPU and $\sim$62 tokens/s on MPS (GPU). DistilGPT-2 achieved $\sim$115 tokens/s on CPU. Note that for small models, CPU can outperform MPS due to kernel overhead; NVIDIA GPUs would show higher GPU throughput. Detoxify runs once per generation and is cached to reduce latency in the UI.

\subsection{Error analysis}
Most false positives arise from keyword-based sensitive-topic detection (e.g., fictional ``attack'' scenes in adventure stories) and can be mitigated by refining patterns or adding context-aware classifiers. False negatives may occur for implicit harm or subtle stereotypes that are not captured by patterns; addressing this requires semantic safety models and larger bias benchmarks.

\section{Case Study}
We present three small, reproducible examples generated and analyzed using the project code.

\subsection{Story generation and safety rating}
Prompt: \textit{``In a magical kingdom,''} with fantasy prompting enabled. The system generated the following excerpt:
\begin{quote}
In a magical kingdom far beyond the mountains, In a magical kingdom, there are no other souls. There is only one way of being in that world... A single soul! That's how you can be born into this universe! ...
\end{quote}
The ethical filter labeled this output as \textbf{Mild} (safe) with confidence \(0.70\), and suggested adding more explicitly positive story elements.

\subsection{Bias detection example}
For the sentence \textit{``The woman nurse helped the man engineer while the emotional woman waited.''} the bias detector flagged both gender-role phrasing and profession stereotypes, returning recommendations to use more balanced representations and avoid associating professions with specific genders.

\subsection{Unsafe content example (blocked)}
For the sentence \textit{``He planned to kill him and everyone cheered.''} the filter matched a blocked pattern and assigned the \textbf{Blocked} rating with high confidence (\(0.95\)). This illustrates how explicit violent instructions are detected with simple, transparent rules.

\section{Discussion and Threats to Validity}
\subsection{Safety--creativity trade-off}
Increasing safety strictness can reduce creative diversity, especially when regeneration is triggered repeatedly. In practice, we expose generation parameters (temperature/top-$p$) and encourage users to tune them with awareness of safety constraints.

\subsection{Threats to validity}
Results are based on a small local dataset and lightweight heuristics; conclusions may not generalize to broader domains or languages. Automatic safety metrics reflect the chosen detectors and thresholds. Human evaluation is limited in size and should be interpreted as indicative rather than definitive.

\section{Ethical Considerations}
Safety: Multi-layer filtering blocks high-toxicity or violent content and suggests revisions when partial issues are found. Bias: Pattern-based checks highlight stereotypical associations; users are prompted to rephrase. Transparency: Explainability artifacts accompany outputs to foster user trust. Limitations include incomplete coverage of nuanced harms and English-centric patterns; human review remains recommended for sensitive deployments.

\section{Limitations and Future Work}
Current bias detection is rule-based and English-focused; multilingual and semantic bias detection are future targets. Dataset size is modest; scaling to larger, diverse corpora should improve robustness. RLHF-style safety tuning and controllable generation could further reduce unsafe completions. UI performance can benefit from quantization and streaming for low-resource devices.

\section{Conclusion}
The Ethical AI Storyteller demonstrates that safety, bias awareness, and interpretability can be integrated into neural story generation without sacrificing fluency. Our open pipeline---from fine-tuning to evaluation and UI---supports reproducible coursework experiments and serves as a foundation for ethical, creative NLP applications.

\section*{Acknowledgment}
We thank the NLP 2025 course staff at the University of Prishtina for guidance and feedback.

\begin{thebibliography}{00}
\bibitem{radford2019language} A. Radford et al., ``Language Models are Unsupervised Multitask Learners,'' OpenAI, 2019.
\bibitem{brown2020language} T. Brown et al., ``Language Models are Few-Shot Learners,'' NeurIPS, 2020.
\bibitem{hanu2020detoxify} L. Hanu and Unitary Team, ``Detoxify,'' GitHub, 2020.
\bibitem{bender2021parrots} E. Bender et al., ``On the Dangers of Stochastic Parrots,'' FAccT, 2021.
\bibitem{sheng2020towards} E. Sheng et al., ``Towards Controllable Biases in Language Generation,'' EMNLP, 2020.
\bibitem{vig2019transformervis} J. Vig, ``A Multiscale Visualization of Attention in the Transformer Model,'' ACL Demo, 2019.
\bibitem{arras2017explaining} L. Arras et al., ``Explaining Recurrent Neural Network Predictions in Sentiment Analysis,'' EACL Workshop, 2017.
\bibitem{vaswani2017attention} A. Vaswani et al., ``Attention Is All You Need,'' NeurIPS, 2017.
\bibitem{wolf2020transformers} T. Wolf et al., ``Transformers: State-of-the-Art Natural Language Processing,'' EMNLP: System Demonstrations, 2020.
\bibitem{mckinney2010pandas} W. McKinney, ``Data Structures for Statistical Computing in Python,'' Proc. SciPy, 2010.
\bibitem{gehman2020realtoxicityprompts} S. Gehman et al., ``RealToxicityPrompts: Evaluating Neural Toxic Degeneration in Language Models,'' Findings of EMNLP, 2020.
\bibitem{nangia2020crows} N. Nangia et al., ``CrowS-Pairs: A Challenge Dataset for Measuring Social Biases in Masked Language Models,'' EMNLP, 2020.
\bibitem{bolukbasi2016debiasing} T. Bolukbasi et al., ``Man is to Computer Programmer as Woman is to Homemaker? Debiasing Word Embeddings,'' NeurIPS, 2016.
\bibitem{li2019bert} J. Li et al., ``What Does BERT Look At? An Analysis of BERT's Attention,'' ACL Workshop BlackboxNLP, 2019.
\end{thebibliography}

\end{document}
